\ifthenelse{\equal{\LanguageOption}{italian}}{%
	\chapter[Introduzione al template: motivazioni e primi passi]{Introduzione al template\\Motivazioni e primi passi}
	\label{cp:introduzione}
	
	{
		\newpage
		\parindent0pt
		\textit{Autore: Antonio José Gregorio Rega} \\
		\textit{Versione Attuale: 1.0.0 (Unibo Edition)} \\
		\textit{Licenza: \LaTeX~Project Public License v1.3c} \\
		\textit{Repository Ufficiale: \href{https://github.com/antoniojoserega/unibo-thesis}{GitHub Repository}}
		
		\vspace{.935em}
		
		Benvenuto nel template \textcolor{unibored}{\textit{Unibo-Thesis}}! Grazie per averlo scelto per la tua tesi o il tuo progetto. Questo template è frutto di un lungo lavoro di ottimizzazione per adattare gli standard accademici dell'Università di Bologna a un sistema moderno e automatizzato. Questo capitolo ti aiuterà a muovere i primi passi. Consulta il \autoref{cp:user-guide} per una guida dettagliata alle funzioni avanzate.
	}
	
	\section{Motivazione}
	Scrivere una tesi è già abbastanza complesso senza dover lottare con i margini di Word o template LaTeX obsoleti e disorganizzati. \textit{Unibo-Thesis} nasce con l'obiettivo di fornire agli studenti dell'Alma Mater uno strumento: \(i)\) esteticamente professionale e conforme ai colori di Ateneo, \(ii)\) estremamente organizzato in cartelle logiche, \(iii)\) automatizzato tramite GitHub Actions per garantire che il PDF si compili sempre correttamente.
	
	\section{Iniziare a scrivere}
	Il modo più semplice per usare questo template è tramite \href{https://www.overleaf.com/}{Overleaf}. Basta caricare il file ZIP del progetto e iniziare a scrivere. Se preferisci lavorare in locale, ti consiglio \href{https://www.tug.org/texlive/}{TeX-Live} o \href{https://miktex.org/}{MikTeX} insieme a un editor come VS Code o TeXstudio.
	
	\begin{block}[tip]
		Il punto di forza di questo template è l'integrazione con GitHub. Ogni volta che fai un \texttt{push} delle tue modifiche, una \textbf{GitHub Action} proverà a compilare il documento. Se il quadratino diventa verde, significa che la tua tesi è tecnicamente perfetta e pronta per la consegna!
	\end{block}
	
	\section{Supporto e contributi}
	Se incontri un bug o hai un suggerimento, apri una \textit{Issue} sul repository GitHub. Se invece questo template ti ha salvato la vita (o almeno la sessione di laurea), puoi supportare il mio lavoro offrendomi un caffè sul mio profilo \href{https://ko-fi.com/antoniojosegregoriorega}{\textbf{Ko-fi}}. Il tuo contributo mi aiuta a mantenere il codice aggiornato per i futuri laureandi!
	
}{%
	\chapter[Introduction to the template: motivation and first steps]{Introduction to the template\\Motivation and first steps}
	\label{cp:introduction}
	
	{
		\newpage
		\parindent0pt
		\textit{Author: Antonio José Gregorio Rega} \\
		\textit{Current Version: 1.0.0 (Unibo Edition)} \\
		\textit{License: \LaTeX~Project Public License v1.3c} \\
		\textit{Official Repository: \href{https://github.com/antoniojoserega/unibo-thesis}{GitHub Repository}}
		
		\vspace{.935em}
		
		Welcome to the \textcolor{unibored}{\textit{Unibo-Thesis}} template! This project is designed to provide University of Bologna students with a modern, automated, and professionally styled LaTeX environment. This chapter introduces its purpose; see \autoref{cp:user-guide} for a detailed guide on how to use its custom features.
	}
	
	\section{Motivation}
	Most academic templates are either too rigid or poorly organized. \textit{Unibo-Thesis} was created to offer: \(i)\) a clean design using official Unibo colors, \(ii)\) a modular file structure, and \(iii)\) built-in automation via GitHub Actions to ensure your document always compiles flawlessly.
	
	\section{Getting started}
	For a hassle-free experience, use \href{https://www.overleaf.com/}{Overleaf} by uploading the project ZIP. If you prefer local development, I recommend \href{https://www.tug.org/texlive/}{TeX-Live} or \href{https://miktex.org/}{MikTeX}. 
	
	\section{Support and Feedback}
	Found a bug? Open an \textit{Issue} on GitHub. If this template saved you hours of formatting stress, please consider supporting my work by buying me a coffee on \href{https://ko-fi.com/your-profile}{\textbf{Ko-fi}}. Your support keeps this project alive for future graduates!
}

\MediaOptionLogicBlank